\section{An actual nuclear return to the local source and its trace}
\label{sec:local-source-return}

The ideal projector in MA-MO-012 cannot be omitted by formal
cyclicity. There is, however, a different possible comparison:
return the completed ideal operator to the actual smooth local
source before taking its trace. We prove that comparison here.
The resulting trace is a local source trace. It is not the trace
of the full narrow zero quotient or of an unconstructed ambient
common-space matrix.

Fix a nontrivial primitive finite-order character \(\eta\), the
structural finite place set \(S\), the actual norm splitting and
local source Fourier maps of MA-CC/MA-CQ/MA-MO-012. Write
\(\mathcal T_\eta\) for the corresponding compact-character
sector of Meyer's \(\mathcal T(C_S)\), with its actual local
coinvariant topology. Use the fixed probability Haar measure
on the compact norm-one factor and center all norm coordinates,
as in MA-CC/MA-CQ. Its global
summation embedding is
\[
 \vartheta_\eta:\mathcal T_\eta\longrightarrow
 I_\eta^c\subset\mathcal E,\qquad
 J_{\eta,a,s}=\overline{I_\eta^c}^{\,H_{\Sigma,a,s}}.
 \tag{MA-MO.280}
\]
This is the original source embedding, not a chosen radial
representative for an ideal class. On the fixed compact sector
the source is a nuclear Fr\'echet space. In the centered norm
coordinate its topology is measured by the Schwartz seminorms
of \(e^{cx}F\) and of the actual dual Fourier vector
\(e^{cx}\mathfrak S_\eta F\), for every \(c>-1/2\).
The dual vector belongs to the conjugate compact sector; finite
Fourier phases and conductor normalizations are retained in
\(\mathfrak S_\eta\).

Let \(R F(x)=F(-x)\). The equivariant local Fourier operator
\(\mathfrak U_\eta=\mathfrak S_{\bar\eta}R\), with sectors
chosen so that it acts on \(\mathcal T_\eta\), has scalar
Mellin multiplier \(m_\eta(z)\). Equivalently this notation
defines \(m_\eta\) by the actual local map; no conductor phase
is removed. The elementary local functional equations give:
\[
 \begin{gathered}
 m_\eta\text{ is holomorphic for }\Re z>-1/2,\\
 m_\eta\text{ and all its derivatives have polynomial growth}\\
 \text{on every fixed closed vertical substrip there},\\
 m_\eta(b+it)^{-1}\text{ and all its derivatives have polynomial growth}\\
 \text{for }-1/2<b<1/2.
 \end{gathered}
 \tag{MA-MO.281}
\]
Here and below polynomial exponents may depend on the fixed
strip and derivative order. These assertions concern the local
multiplier only. For the rational even archimedean sector, for
example, with \(\Gamma_{\mathbb R}(w)=\pi^{-w/2}\Gamma(w/2)\),
\[
       m(z)=\frac{\Gamma_{\mathbb R}(1/2+z)}
                       {\Gamma_{\mathbb R}(1/2-z)}.
 \tag{MA-MO.282}
\]
At a real place the parity shifts both gamma arguments by zero
or one; a complex finite-order sector has its corresponding
complex gamma quotient. Up to the nonzero exponential monomial
from the different and the chosen additive character, an
unramified finite factor is a ratio
\((1-\omega q^{z-1/2})/(1-\omega' q^{-z-1/2})\), with unit
modulus constants in the numerator and denominator; ramified
factors are nonzero exponential monomials times unit constants.
Thus their possible poles lie on \(\Re z=-1/2\), and their
possible positive-strip zeros are on \(\Re z=1/2\).
Stirling's formula proves the gamma bounds, including derivatives;
the finite factors and their derivatives are bounded on each
fixed admissible input line. This verifies (MA-MO.281) without
using any global zero location.

Let \(\phi=1\) sufficiently far left and \(\phi=0\)
sufficiently far right, as before, and put
\[
 \begin{aligned}
 \phi^\vee(x)&=\phi(-x),&
 \psi(x)&=1-\phi(x)-\phi(-x),\\
 T_{\eta,\phi}&=M_\phi-
       \mathfrak U_\eta M_\phi\mathfrak U_\eta^{-1},&
 \mathcal L_\eta&=-T_{\eta,\phi}
        +\mathfrak U_\eta M_\psi\mathfrak U_\eta^{-1}.
 \end{aligned}
 \tag{MA-MO.283}
\]
The last expression is exactly
\(1-M_\phi-\mathfrak S_{\bar\eta}M_\phi\mathfrak S_\eta\)
on the original source, by inversion of the cutoff. In the
rational even additive coordinates it is the \(L\) of
MA-MO-008. It equals \(-T_{\eta,\phi}\) for a symmetric
cutoff; otherwise the compact residual \(\psi\) is retained.
The actual two Fourier maps conjugate \(\mathcal L_\eta\)
to \(\mathcal L_{\bar\eta}\). No printed Fourier
anticommutation identity is needed.

\paragraph{MA-MO-017 (nuclear local-source return and evaluated ideal trace).}
Assume either analytic regime of MA-MO-012: generic widths with
compact Gevrey \(g,\phi'\), or the stated fixed horizontal zero
gap with arbitrary compact smooth \(g,\phi'\). For every
\(s\ge0\) there is a continuous nuclear map
\[
 \begin{aligned}
 A_{\eta,g}:J_{\eta,a,s}&\longrightarrow\mathcal T_\eta,\\
 A_{\eta,g}\vartheta_\eta F&=C_g\mathcal L_\eta F,\\
 \vartheta_\eta A_{\eta,g}&=B_{\eta,g},
 \end{aligned}
 \tag{MA-MO.284}
\]
where \(B_{\eta,g}\) is the actual completed-ideal endomorphism
of MA-MO-012. Thus its trace has the exact evaluation
\[
 \boxed{\displaystyle
 \operatorname{Tr}_{J_{\eta,a,s}}B_{\eta,g}
   =\frac1{2\pi}\int_{\mathbb R}
          \mathcal Mg(it)\frac{m_\eta'(it)}{m_\eta(it)}\,dt
       +g(0)\int_{\mathbb R}\psi(x)\,dx .}
 \tag{MA-MO.285}
\]
The integral is absolutely convergent. In particular the extra
cutoff term vanishes for the requested \(g(0)=0\) class.
The answer is independent of \(a,s\) within the explicit
regimes. In the notation of the actual local Weil trace,
the first term is \(-W^{\rm loc}_{\eta,S}(g)\), with
all finite-place and different/conductor factors in \(m_\eta\).

The projected trace in (MA-MO.190) therefore equals the right
side of (MA-MO.285). This evaluates that projected trace by a
genuine nuclear factorization; it does not assert
\(\Pi_J=1\) or equality with the unrestricted ambient trace.

There is also a rational principal-sector statement. Let
\(\mathcal W_+\) be the polar-retaining source completion of
MA-MO-016 and \(\jmath_+:\mathscr S_{\rm even}(\mathbb R)
\to\mathcal W_+\) its original theta-pair embedding. The
actual source endomorphism \(C_gL\) has a continuous nuclear
return
\[
 A_{+,g}:\mathcal W_+\longrightarrow
                    \mathscr S_{\rm even}(\mathbb R),\qquad
 A_{+,g}\jmath_+f=\mathcal D_gLf.
 \tag{MA-MO.286}
\]
Its resulting endomorphism \(\jmath_+A_{+,g}\) has trace
(MA-MO.285) with (MA-MO.282). This is distinct from the full
source column \(C_g(P_+-P_-)|_{\mathcal W_+}\), whose two
outgoing polar columns are explicitly retained in MA-MO-016.

\begin{proof}
\emph{An exact kernel between different exponential weights.}
Write \(V_b=\mathcal F_+M_{e^{bx}}\), with unitary Fourier
normalization, and set
\[
 \begin{gathered}
 z=c+it,\quad w=b+iu,\quad G(z)=\mathcal Mg(z),\\
 P(z-w)=\int e^{(z-w)x}\phi'(x)\,dx,\qquad
 \Psi(z-w)=\int e^{(z-w)x}\psi(x)\,dx.
 \end{gathered}
 \tag{MA-MO.287}
\]
For every \(c>-1/2\) and fixed \(-1/2<b<1/2\), the
kernel of \(V_c C_g\mathcal L_\eta V_b^{-1}\) on its
original smooth source is
\[
 \ell_{c,b}(t,u)=\frac{G(z)}{2\pi}
 \left[
 -\frac{P(z-w)}{z-w}
              \left(\frac{m_\eta(z)}{m_\eta(w)}-1\right)
       +\frac{m_\eta(z)}{m_\eta(w)}\Psi(z-w)
 \right].
 \tag{MA-MO.288}
\]
At \(z=w\) the divided difference has its removable value.

Here is a direct derivation, so that no inverse local multiplier
is moved across one of its positive zeros. If \(h\) and its
two exponential weights are Schwartz, the kernel of
\([M_\phi,C_h]\), between these weighted lines, has double
Fourier transform
\[
 -\frac{P(z-w)}{2\pi(z-w)}
              \bigl(\mathcal Mh(w)-\mathcal Mh(z)\bigr).
 \tag{MA-MO.289}
\]
Indeed, putting \(r=x-y\), integration in \(y\) uses the
compactly supported difference \(\phi(y+r)-\phi(y)\).
An integration by parts gives
\[
 \int e^{(z-w)y}\{\phi(y+r)-\phi(y)\}\,dy
       =-\frac{e^{-(z-w)r}-1}{z-w}P(z-w).
 \tag{MA-MO.290}
\]
This identity is valid for either sign of \(c-b\); it is not
an unjustified transform of the nonintegrable step itself.
At equality use its continuous limit. The decay hypotheses on
\(h\) justify the remaining integration and differentiation.

By (MA-MO.281), the Mellin transform of
\(\mathfrak U_\eta g\) is \(m_\eta(z)G(z)\), rapidly
decreasing on every fixed positive finite-\(L\) half-plane
line. Contour shifting, or the actual local Fourier estimate,
places this function in all the required weighted Schwartz
spaces. Now use the exact operator identity
\[
 C_gT_{\eta,\phi}
     =-[M_\phi,C_g]
       +[M_\phi,C_{\mathfrak U_\eta g}]
                    \mathfrak U_\eta^{-1}.
 \tag{MA-MO.291}
\]
Substitution of (MA-MO.289) cancels the two terms containing
\(G(w)\). It gives the negative of the first summand in
(MA-MO.288). The compact multiplier term in (MA-MO.283)
has kernel \(G(z)m_\eta(z)\Psi(z-w)/(2\pi m_\eta(w))\).
This proves (MA-MO.288), with the inverse evaluated only on
the fixed input line \(b\).

\emph{Composition with the genuine global inverse.}
Use the actual finite Euler multiplier of MA-MO-012:
\[
 Z_\eta(z)=E_{\eta,S}(1/2+z)L_\eta(1/2+z),
       \qquad j_{\eta,b}(u)=Z_\eta(b+iu)^{-1}.
 \tag{MA-MO.292}
\]
For every smooth ideal vector \(y=\vartheta_\eta F\),
the unconditional division result (MA-MO.173--178) identifies
its original inverse with \(V_bF=j_{\eta,b}V_by\).
It also supplies the corresponding identity for the actual
dual Fourier vector and \(Ry\). Thus the kernel giving the
first source component of the proposed return is exactly
\[
                  \ell_{c,b}(t,u)j_{\eta,b}(u).
 \tag{MA-MO.293}
\]
There is no assertion that \(j_{\eta,b}\) separately acts
on every Schwartz vector.

For clarity, (MA-MO.293) is a Schwartz function of \((t,u)\)
for every \(c>-1/2\). In the gap regime all derivatives of
\(j_{\eta,b}\) are polynomially bounded; (MA-MO.281), the
rapid transforms \(G\), \(P\), \(\Psi\), and the divided
difference prove the assertion. In the generic regime their
derivatives are bounded by \(\exp(C\log^2(2+|u|))\).
The compact residual \(\psi\) has the same Gevrey order as
\(\phi'\), since its derivatives are sums of derivatives of
that fixed step and its reflection. Compact Gevrey data give decay
\[
 C_N(1+|t|+|t-u|)^{N}
       \exp\{-d(|t|^{1/q}+|t-u|^{1/q})\}
 \tag{MA-MO.294}
\]
for suitable constants in each derivative estimate. The
exponential defeats that logarithmic-square growth since
\(|u|\le |t|+|t-u|\). The divided difference creates no
singularity: write
\((m(z)-m(w))/(z-w)=\int_0^1m'(w+v(z-w))\,dv\).
The segment stays in a fixed positive finite-\(L\) half-plane
strip. All remaining factors have the polynomial bounds in
(MA-MO.281). These estimates hold in every derivative order
and locally uniformly in \(c\).

Every such kernel maps \(L^2(du)\) continuously into
\(\mathscr S(dt)\), by Cauchy--Schwarz applied to its
weighted derivatives. The input map \(y\mapsto V_by\)
is bounded from \(H_{\Sigma,a,s}\) to \(L^2\), for
\(b=a\) or \(-a\), since \(s\ge0\). Consequently all
weighted source seminorms of the first output are bounded by
the declared Hilbert norm of \(y\).

Apply exactly the same construction in the dual compact
character sector to \(Ry\), with \(g^\vee(x)=g(-x)\).
The actual source Fourier identities and (MA-MO.283) give
\[
 \mathfrak S_\eta C_g\mathcal L_\eta F
       =C_{g^\vee}\mathcal L_{\bar\eta}
                            \mathfrak S_\eta F.
 \tag{MA-MO.295}
\]
Thus the second family of kernels controls every seminorm of
the dual Fourier output as well. This step retains the dual
character and the actual finite Fourier transform.

Take a sequence from the smooth ideal tending to \(y\in J\).
The two outputs converge in their weighted Schwartz families.
At the central local \(L^2\) weight the source Fourier map is
unitary, so their limits remain Fourier transforms of each
other. The source topology described above therefore puts the
limit in \(\mathcal T_\eta\), continuously. It is independent
of the approximating sequence. This constructs \(A_{\eta,g}\)
and proves its first identity on the complete smooth source.

\emph{Nuclearity into the actual source topology.}
We spell out the functional-analytic point used here. The fixed
compact-character source is nuclear Fr\'echet, with a cofinal
increasing family of Hilbertian norms \(q_n\), each dominating
the injective central local \(L^2\) norm, whose successive
completion maps may be chosen trace class. Thus none of the
finite-dimensional source images below has a seminorm-null
vector. A
continuous map from a Hilbert space into this source is trace
class in each \(q_n\): factor it through a stronger seminorm
and the corresponding trace-class completion map. Choose
finite-rank orthogonal domain projections \(P_N\) increasing
strongly to the identity. For each \(n\), the \(q_n\)
trace norm of \(A(1-P_N)\) tends to zero. Choose an increasing
subsequence \(N_k\) making these tails at most \(2^{-k}\)
for every \(n\le k\). Decompose each finite-rank block
\(A(P_{N_{k+1}}-P_{N_k})\) by its singular-value expansion
in \(q_k\). Its vectors lie in the actual source, its
coefficient functionals have norm one, and the sum of its
singular values is at most \(2^{-k}\). For each fixed
\(q_n\), the resulting normalized source vectors are bounded
once \(k\ge n\); the finitely many earlier blocks are bounded
as well. Together with the first finite block this is a
summable nuclear series with a bounded set of actual source
vectors. This proves the required nuclearity, rather than
only separate one-sided Hilbert trace classes.

The source's nuclear Fr\'echet property and Hilbertian
description are the standard fixed-sector consequence of
Meyer's split extension (24), Proposition 4.4 and the
weighted Schwartz topology; the same bounded-map argument is
used in Proposition 4.6. Alternatively, in the rational
archimedean sector one may use the even Hermite basis of
\(\mathscr S(\mathbb R)\) directly. Its coefficient norms
decay faster than every power by the source seminorm estimates
just proved and give the same nuclear series.

\emph{Actual factorization and the trace.}
The global source embedding (MA-MO.280) is continuous into
\(J_{\eta,a,s}\): its two Poisson-related tails give every
exponential seminorm for a nontrivial character, as in
MA-MO-012. The first source identity in (MA-MO.284) and the
literal common-sum formula imply
\(\vartheta_\eta A_{\eta,g}=B_{\eta,g}\) on a dense smooth
ideal. Both sides are continuous, so the identity holds on
the Hilbert completion. The reverse composition is the
original local operator \(C_g\mathcal L_\eta\) on
\(\mathcal T_\eta\). Nuclear trace cyclicity for these
actual maps therefore gives
\[
 \operatorname{Tr}_J(\vartheta_\eta A_{\eta,g})
       =\operatorname{Tr}_{\mathcal T_\eta}
                    (A_{\eta,g}\vartheta_\eta)
       =\operatorname{Tr}_{\mathcal T_\eta}
                              C_g\mathcal L_\eta.
 \tag{MA-MO.296}
\]
Concretely both traces are the same absolutely summable
pairing series from the preceding nuclear expansion. The
smooth source has the approximation property from its split
weighted-Schwartz model, so this is its canonical nuclear
trace. Meyer's Proposition 4.6 and Corollary 4.7 identify it
with the local central \(L^2\) trace; the compact residual
term has the same comparison by the elementary kernel lemma.
Equivalently, repeating (MA-MO.288--295) with the global
inverse factor replaced by one proves the local central
\(L^2\)-to-source nuclear factorization directly. This
also avoids relying on the printed Fourier-sign shorthand
in the source proof of Proposition 4.6.

Its normalization may also be read directly from (MA-MO.288)
without the global inverse, at \(c=b=0\). That kernel is
Schwartz and hence trace class. Since \(P(0)=-1\), its
diagonal is
\[
 \ell_{0,0}(t,t)=\frac{G(it)}{2\pi}
       \left(\frac{m_\eta'(it)}{m_\eta(it)}
                                      +\int\psi\right).
 \tag{MA-MO.297}
\]
The Schwartz-kernel trace equals its diagonal integral, and
\((2\pi)^{-1}\int G(it)\,dt=g(0)\). Formula (MA-MO.285)
follows, including its sign and compact-residual term.
The logarithmic derivative of the local gamma/finite factors
has at most polynomial growth, so the displayed integral is
absolutely convergent. This is also the local Weil trace of
Meyer's equation (20), with the sign from (MA-MO.283).

\emph{The rational polar source.}
On \(I=\Theta\mathscr U_0\), the same construction uses
the actual inverse identities in MA-MO-010. The paired
Fourier estimates put its output in the actual even additive
Schwartz space. This can be verified without relying on a
radial description alone: taking \(c=-1/2+\varepsilon\)
gives local integrability at additive zero, all positive
\(c\) give rapid decay at infinity, and the paired Fourier
estimates give integrability of all Fourier moments. Fourier
inversion then gives smooth bounded derivatives at zero and
the ordinary Schwartz estimates. This verifies the actual
additive target.

For \((y,c_0,c_1)\in J\oplus\mathbb C^2=\mathcal W_+\),
using the two actual source sections of MA-MO-016, define
\[
 A_{+,g}(y,c_0,c_1)
       =A_{0,g}y+\sum_{j=0}^1c_j\mathcal D_gLf_j.
 \tag{MA-MO.298}
\]
The extra columns are actual Schwartz vectors and have finite
rank. The result is continuous and nuclear, and on the dense
original positive source it is exactly \(\mathcal D_gL\).
The theta-pair inclusion into \(\mathcal W_+\) is continuous
by the seminorm estimates in MA-MO-016. Applying the same
nuclear pairing argument proves (MA-MO.286) and its trace.
On the polar sections \(i_+f_j\), however,
\[
 (P_+-P_-)i_+f_j
       =i_+Lf_j+(i_+f_j-P_-i_+f_j).
 \tag{MA-MO.299}
\]
The latter finite columns are precisely why this source
endomorphism is not silently identified with the full raw
matrix defect. This completes the proof.
\end{proof}

\paragraph{Primary locators and the remaining arithmetic gate.}
The supplied Meyer (2005) source is used at Lemma 4.3
(author-manuscript statement p.~25, proof p.~26),
Proposition 4.4 (pp.~26--27),
the source topology and split extension (24) (pp.~27--28),
Proposition 4.6 and Corollary 4.7 (p.~28), and the
local trace equation (20) (p.~22; (19) is its normalized-character
version). Their exact TeX labels are
\texttt{lem:mother\_of\_nuclearity},
\texttt{pro:Fourier\_estimate},
\texttt{eq:Twist\_convariants},
\texttt{pro:T\_on\_Twist}, \texttt{cor:T\_on\_Twist}, and
\texttt{eq:local\_trace\_formula\_psi}.
The newly needed assertion is the return after the possibly
unbounded global inverse; it is (MA-MO.288--295), not a direct
application of the local theorem to that inverse.

The elementary local factors in (MA-MO.281--282) are Tate
(1950), Theorem 2.4.1, thesis p.~(2.11), with the explicit
tables in \S2.5: real p.~(2.14), complex p.~(2.16), and
nonarchimedean p.~(2.19). These are PDF pages 18, 21, 23,
and 26 of the original-thesis scan. Those tables are used for
the gamma/Euler forms and their analytic domains; the actual
Fourier convention, different, conductor and phases here are
fixed by the source definition of \(m_\eta\), not transferred
from an inconsistent Fourier-sign convention.

\smallskip
\noindent Tate, J. (1950). \emph{Fourier analysis in number
fields and Hecke's zeta-functions} [Doctoral dissertation,
Princeton University]. Original-thesis scan,
\url{https://sites.math.rutgers.edu/~alexk/2023S572/Tate1950.pdf}.

No zero-location premise enters the local multiplier or
cross-weight construction. Global reciprocal bounds use the
explicit generic/Gevrey or gap/smooth hypotheses. For rational
and Dirichlet sectors A supplies the gap at every positive
narrow width; using it is recorded as zero-location input.
For the non-induced character, M or its own RH supplies that
gap, while the generic/Gevrey alternative remains unconditional.
The full quotient trace and the FT002 discarded-divisor
condition are unchanged. The present theorem closes the
specific evaluation of the completed positive-ideal source
trace, not the independent full narrow identity.

In the inherited cote~1 framework, scans~079--080 and~111--112
through ALPOGE\_BRIDGES BR-TR-01/02, (MA-MO.284) is an actual
nuclear factorization with compatible pairings. It therefore
justifies a trace comparison and its associated Fredholm
endomorphism determinant. The finite-order character remains
the specified rank-one finite-monodromy local system; no
positive faithful full-jet realization or new arithmetic weight
is inferred from this nuclear comparison.
